% !TeX root = ../navier-stokes-manuscript.tex
% ============================================================
% 1.5 Class Membership and Participation
% ============================================================
\subsection{Class Membership and Participation}\label{sub:ClassMembershipAndParticipation}

Silver ends with a practical difficulty.
Every proposed loss of smoothness has to be tested against the same original fluid, so we need an exact way to say whether the proposed object still belongs to that fluid at all.

Fix the initial datum \(u_0\) on \(\mathbb R^3\) and one viscosity \(\nu>0\).
Let \((u,p)\) be the classical solution generated by those choices on its maximal interval \([0,T_*)\).
Let \(Q=(\widetilde u,\widetilde p)\) be a proposed extension of this particular preterminal solution to \([0,T_*+\varepsilon)\), for some \(\varepsilon>0\).
Every membership test below concerns what this fixed evolution can become after \(T_*\); weak solutions constructed from unrelated data are outside the question.

Write \(\operatorname{Member}(Q)\) when the proposed extension agrees with \((u,p)\) before \(T_*\) and continues to obey the requirements fixed at the beginning:
\[
  \partial_t\widetilde u
  +(\widetilde u\cdot\nabla)\widetilde u
  +\nabla\widetilde p
  =
  \nu\Delta\widetilde u,
  \qquad
  \nabla\cdot\widetilde u=0,
\]
with the same datum, viscosity, decay, energy condition, and pressure determined by the same velocity field.
For a nonsmooth proposal, the equations are read in the sense of distributions.
The energy condition means
\[
  \widetilde u
  \in
  L^\infty_{\mathrm{loc}}\bigl([0,T_*+\varepsilon);L^2(\mathbb R^3)\bigr)
  \cap
  L^2_{\mathrm{loc}}\bigl([0,T_*+\varepsilon);H^1(\mathbb R^3)\bigr),
\]
together with the Leray energy inequality.
Smoothness remains a separate conclusion to be proved.

Class exit is simply the complement
\[
  \operatorname{Exit}(Q)
  :=
  \neg\operatorname{Member}(Q).
\]
Gold and Silver face the same boundary from opposite directions.
Gold seeks
\[
  \operatorname{Member}(Q)
  \Longrightarrow
  \operatorname{Smooth}(Q),
\]
while Silver seeks the contrapositive
\[
  \neg\operatorname{Smooth}(Q)
  \Longrightarrow
  \operatorname{Exit}(Q).
\]
The two implications have the same logical content.
Among the proposed extensions of this one history, put
\[
  \mathcal M
  :=
  \{Q:\operatorname{Member}(Q)\},
  \qquad
  \mathcal S
  :=
  \{Q:\operatorname{Smooth}(Q)\}.
\]
Gold and Silver make the same set-theoretic claim:
\[
  \mathcal M\subseteq\mathcal S,
  \qquad\text{or equivalently}\qquad
  \mathcal M\cap\mathcal S^{\mathsf c}=\varnothing.
\]
Gold begins with an admitted member and tries to place it inside the smooth set.
Silver begins with a proposed nonsmooth extension and tries to place it outside the membership set.
That equivalence fixes the direction of the Silver argument; the class-exit implication itself still has to be proved.
A large derivative, an unusual geometry, or a suggestive limiting picture gives evidence only after it identifies a requirement that the proposed extension can no longer satisfy.

\divider{What Participation Means}

To see what those requirements mean physically, follow one material parcel through the field.
If \(X(a,t)\) is the trajectory that begins at the material label \(a\), then
\[
  \frac{d}{dt}X(a,t)=u(t,X(a,t)).
\]
The parcel is carried into a new neighbourhood by the velocity itself.
Its velocity along that path changes according to
\[
  \frac{d}{dt}u(t,X(a,t))
  =
  D_tu(t,X(a,t))
  =
  -\nabla p(t,X(a,t))
  +\nu\Delta u(t,X(a,t)),
\]
where \(D_t=\partial_t+u\cdot\nabla\) is the material derivative.

The parcel moves into the velocity field around it, the nearby curvature of that field supplies the viscous part of its acceleration, and the pressure gradient comes from the incompressibility constraint imposed across the whole fluid.
As the parcel moves, the surrounding velocity and the geometry of its material neighbourhood change.
As the surrounding velocity changes, the local viscous response and the global pressure response are both determined again by the changed field.
The parcel never acquires an independent equation of motion detached from that field.

This continuing dependence is what we mean by \emph{participation}.
A parcel participates when its motion remains part of the same coupled evolution: it is transported by \(u\), it changes velocity through the pressure and viscous terms generated by that \(u\), and it remains inside the divergence-free motion of the whole fluid.
A small region participates when the same relation holds throughout the region and agrees with the surrounding field along its boundary.

Participation allows perfectly ordinary zeroes.
A parcel may be instantaneously motionless even while it is accelerating.
Its material acceleration vanishes only when the pressure and viscous terms balance.
The zero solution has every derivative equal to zero and still satisfies the complete relation.
The decisive issue is whether each value continues to agree with the coupled field that determines it.

Couette flow gives one part of this relation a simple physical picture.
Place fluid between the plates \(x_2=0\) and \(x_2=h\), hold the lower plate fixed, and move the upper plate with constant tangential velocity \(\gamma h e_1\).
The steady interior profile is
\[
  u(x)=(\gamma x_2,0,0),
  \qquad
  p=\text{constant}.
\]
Its shear per unit density is
\[
  \tau_{12}=\nu\partial_{x_2}u_1=\nu\gamma.
\]
The nonzero shear expresses the transfer of tangential momentum from one layer to the next.

The caveat matters.
This linear profile also satisfies
\[
  \Delta u=0,
  \qquad
  (u\cdot\nabla)u=0.
\]
The viscous stresses on the two sides of an interior layer balance, so the bulk viscous force and the material acceleration are both zero there.
The moving plate continuously supplies the momentum that maintains the profile.
Couette flow isolates viscous shear between neighbouring layers.
Its moving boundaries keep it separate from an unforced whole-space evolution, and its zero Laplacian shows why a nonzero gradient alone does not force a nonzero acceleration.

In a general flow the local balance changes from one parcel to the next, transport changes the geometry and relative spacing of material neighbours, and pressure keeps the whole motion incompressible at the same time.
An admitted smooth continuation participates in this complete relation.
A concrete failure of one of its terms can witness class exit, while loss of continuity or loss of finite control requires its own endpoint test.
